医学图像分割是计算机辅助诊断中的基础任务。在临床实践中，医生通过大致勾画目标区域的边界框（bounding box）来指示感兴趣的器官或病灶，模型据此给出精确的二值分割掩码——这一交互式分割范式兼顾了自动化效率与临床可控性。MedSAM 作为 Segment Anything Model（SAM）在医学领域的迁移模型，为上述提示式医学分割提供了统一基座，但在腹部多器官 CT 分割任务中仍存在两方面不足：box prompt 驱动的二值分割训练易受前景/背景失衡与困难像素学习不充分的影响；ViT 编码器侧重全局语义建模，对局部边界与细粒度结构的表征能力有限。本文以 MedSAM 为基座、以腹部多器官 CT 为统一验证场景，聚焦上述两个问题开展系统优化研究。

针对训练信号分配问题，本文设计了联合平衡损失框架 Balance Loss，从类别间重平衡与类别内困难像素加权两个层面重构像素级监督信号，并结合两阶段训练策略缓解早期困难样本挖掘噪声。针对局部边界表征不足，本文在固定损失主干后比较多种微调与结构增强路线，最终在 MedSAM 的特征桥接位置引入轻量级多尺度局部适配器（MSL-Adapter），通过双路径深度可分离卷积与残差注入补充局部空间信息与边界敏感响应。冻结主干 LoRA 微调等负向结果亦作为路线筛选的排除性证据予以保留，用以阐明不同策略的适用边界。

实验在统一评估协议下开展：联合 FLARE22 与 AMOS22 CT 子集共 60 个训练病例，在 15 个无重叠测试病例上以 GT box prompt 条件进行病例级评估。结果表明，Balance Loss 可稳定改善基线在困难器官及边界指标上的表现；局部适配路线进一步带来边界质量增益。配对 Wilcoxon 符号秩检验经 Bonferroni 校正后确认 Baseline 至最终方案的差异具有统计显著性。与已有方法的横向对比显示，本文方案在相同提示条件下优于 SAMed 和 Medical SAM Adapter 等同类方法；与全自动方法 nnU-Net 的跨范式参考性对比中，由于 GT box prompt 为提示式方法提供了全自动方法所不具备的目标区域先验，两类方法解决的任务本质不同，数值不宜做严格横向比较。上述结果表明，"训练信号重构—局部边界补强"的两步优化路径能够有效提升 MedSAM 在腹部多器官 CT 分割任务中的性能。
